% =====================================================================
%  APÊNDICE C — código do robô de médias móveis (02/11/2026, pedido do autor:
%  "apêndice C do modelo replicado do almeida ... a forma mais crua que
%  fizemos, e especificar o que for necessário pra indicar que é uma réplica").
%  O arquivo listado é gerado pelo Makefile a partir de
%  robos/MediasMoveis2/MediasMoveis2.mq5 por ferramentas/scripts/codigo_apendice.sh,
%  com os mesmos marcadores dos Apêndices A e B: saem comentários, a avaliação
%  no testador (Selic, proventos, métricas), StopLoss/TakeProfit, a escolha do
%  método, do preço e do deslocamento das médias (fixos em média simples sobre
%  o fechamento, sem deslocamento) e a opção de operar a cada tick.
%  Os marcadores são só comentários: o EA sem comentários ficou idêntico ao de
%  antes (diff, 02/11/2026). A listagem compila sozinha no prefixo novo
%  (0 erros, 0 avisos, 02/11/2026).
%  Procedência do texto: Almeida (2016), bibliografia/pdfs/Almeida-2016.pdf:
%  seis bibliotecas externas, p. 35; Anexo 1 (EA_TG_2016.mq5), p. 69-77, com
%  RapidaMM_Periodo = 10, DevagarMM_Periodo = 20, método 0 (simples),
%  PRICE_CLOSE, StopLoss = 0 e TakeProfit = 0, compra quando
%  "RapidaMM.Main(Barra_Movimento) > DevagarMM.Main(Barra_Movimento) &&
%  PositionType() != POSITION_TYPE_BUY" (p. 75), gestão de dinheiro por
%  MoneyManagement(), de GerenciaDinheiro.mqh, que não está no anexo
%  (dados/FONTES.md, item 7). Entrada por estado × por cruzamento muda só a
%  primeira operação de cada teste: dados/resultados/replica_almeida_v2_2026-10-21/
%  LEIAME.md, seção 1.
% =====================================================================
\chapter{Código do robô de médias móveis}

Este apêndice apresenta o código do robô de médias móveis descrito no
Capítulo~II, que reproduz a estratégia de \citeonline{almeida2016}. Os
parâmetros das médias mantêm os nomes do anexo original,
\texttt{RapidaMM\_\allowbreak Periodo} e \texttt{DevagarMM\_\allowbreak Periodo}.
Como nos Apêndices~A e~B, foram retirados os comentários, as rotinas de medição
dos resultados e as opções que não foram usadas nos experimentos, como os
limites de perda e de realização de lucro, que o autor deixou com valor zero nos
testes.
% 06/11/2026, pedido do autor: o porquê da reescrita e as diferenças em relação ao
% original foram para o Cap. II (Implementação dos modelos). Texto anterior em
% ia/rascunhos/antes_enxugar_apendice_c_2026-11-06/.

\lstinputlisting[style=mql5]{codigo/MediasMoveis2.mq5}
